% Volume VI: Machine Learning Mathematics
% Continuity note for Chapter 34.
% Previous dependencies: Chapters 13-15 provide joint, marginal, conditional distributions and Bayes' theorem; Chapter 17 provides likelihood and MLE; Chapter 18 provides KL divergence and cross-entropy; Chapters 31-33 provide statistical learning, likelihood losses, regularization, kernels, and Gaussian conditioning.
% New durable concepts: probabilistic graphical models, directed acyclic graphs, undirected factorization, factor graphs, conditional independence, Markov blankets, latent variables, generative direction, inference direction, marginal likelihood, exact posterior computation, sequential Bayes, mixture models, factor analysis intuition, expectation-maximization, responsibilities, ELBO, Jensen lower bound, Laplace approximation, variational inference, Monte Carlo, message passing, particle filtering, and approximate inference caveats.
% Forward hooks: Chapter 41 uses latent-variable inference in VAEs, flows, diffusion, and energy-based models; Chapter 46 uses latent and observed variables in neural GLMs and state-space spike models; Chapter 48 uses latent dynamics; Chapter 50 reuses sequential Bayes and prediction-update loops; Chapter 53 integrates Bayes, predictive coding, and stochastic processes.
% Notation commitments: Observed variables are x or X, latent variables are z or Z, parameters are theta, directed graphical models factor as p(x_1,\ldots,x_d)=prod_j p(x_j|x_{\operatorname{pa}(j)}), undirected models factor as p(x)=Z^{-1} prod_c psi_c(x_c), latent generative models factor as p_theta(x,z)=p_theta(z)p_theta(x|z), posterior inference computes p_theta(z|x), responsibilities are r_{ik}, and variational approximations are q(z).
% Figure plan: no missing images are included. Proposed ImageGen prompts are recorded in imagegen-figures/ch34-figure-metadata.md.
\section{Chapter 34: Probabilistic Graphical Models and Latent Variables}
\label{ch:graphical-models-latent-variables}

Chapter 33 put distributions on functions. Chapter 34 changes the object of probability again: many models describe observed data by introducing hidden causes. A stimulus category may generate a neural response. A discrete state may generate a sequence of spikes. A latent position on a manifold may generate a high-dimensional neural population vector. A topic may generate words in a document. The model runs from latent cause to observation; inference runs backward from observation to latent cause.

The central problem is to organize joint distributions so that structure is visible and computation is possible. Graphs express factorization and conditional independence. Latent-variable models express hidden causes and marginal likelihoods. EM alternates between inferring hidden variables and fitting parameters. Approximate inference replaces exact posteriors when summing or integrating over all latent possibilities is too expensive.

\paragraph{Chapter problem.}
How do graphical structure and latent variables let us build generative probabilistic models, compute posteriors when possible, and approximate them when exact inference is intractable?

\paragraph{Section map.}
Section 34.1 defines directed and undirected graphical models, factorization, and conditional independence. Section 34.2 builds latent-variable generative models and computes exact posteriors in small examples, including sequential Bayes. Section 34.3 derives expectation-maximization as alternating posterior inference and parameter fitting. Section 34.4 surveys approximate inference through Laplace, variational methods, Monte Carlo, message passing, and particle filtering.

\subsection{34.1 Graphical models}

\paragraph{Guiding question.}
How can a graph encode the factorization and conditional-independence structure of a joint probability distribution?

\paragraph{Beginner-facing intuition.}
A joint distribution over many variables can be too large to write as one table. A graphical model uses nodes for random variables and edges for direct probabilistic relationships. The graph does not replace probability; it organizes it. A good graph says which local conditional distributions or local factors multiply together to form the joint distribution.

The main benefit is not the picture. It is the reduced bookkeeping. If a variable depends directly only on a few parents, we do not need a separate parameter for every configuration of all other variables.

\paragraph{Formal definitions.}
A directed graphical model uses a directed acyclic graph, or DAG. Its nodes are random variables \(X_1,\ldots,X_d\). If \(\operatorname{pa}(j)\) is the set of parents of node \(j\), the DAG factorization is
\[
\boxed{
p(x_1,\ldots,x_d)
=\prod_{j=1}^d p\bigl(x_j\mid x_{\operatorname{pa}(j)}\bigr).
}
\]
The graph is acyclic so variables can be ordered from causes or predecessors toward descendants.

An undirected graphical model uses nonnegative potential functions on cliques or factors:
\[
\boxed{
p(x)=\frac{1}{Z}\prod_{c\in\mathcal C}\psi_c(x_c),
}
\]
where \(x_c\) is the subset of variables in clique \(c\), \(\psi_c(x_c)\geq0\), and
\[
Z=\sum_x \prod_{c\in\mathcal C}\psi_c(x_c)
\]
or the corresponding integral normalizes the distribution. The constant \(Z\) can be hard to compute.

A factor graph is a bipartite graph with variable nodes and factor nodes. It explicitly represents a factorization
\[
p(x_1,\ldots,x_d)\propto\prod_a \phi_a(x_{\partial a}),
\]
where \(\partial a\) are the variables touching factor \(a\).

\paragraph{Core mechanism: conditional independence from a simple DAG.}
Consider the directed graph
\[
Z\to X,\qquad Z\to Y.
\]
The DAG factorization is
\[
\boxed{
p(z,x,y)=p(z)p(x\mid z)p(y\mid z).
}
\]
This factorization implies
\[
X\perp Y\mid Z.
\]
Indeed,
\[
p(x,y\mid z)
=\frac{p(z,x,y)}{p(z)}
=p(x\mid z)p(y\mid z).
\]
Thus \(X\) and \(Y\) may be dependent marginally because they share the hidden or observed common cause \(Z\), but they become independent after conditioning on \(Z\). This is the core distinction between marginal dependence and conditional independence.

\paragraph{Concrete example.}
Let \(Z\in\{0,1\}\) be a stimulus state with \(P(Z=1)=0.4\). Let \(X\) and \(Y\) be two binary neural response indicators. Suppose
\[
P(X=1\mid Z=0)=0.1,\quad P(X=1\mid Z=1)=0.8,
\]
and
\[
P(Y=1\mid Z=0)=0.2,\quad P(Y=1\mid Z=1)=0.7.
\]
Then
\[
\begin{aligned}
P(Z=1,X=1,Y=0)
&=P(Z=1)P(X=1\mid Z=1)P(Y=0\mid Z=1)\\
&=0.4\cdot0.8\cdot0.3=0.096.
\end{aligned}
\]
The graph tells us to multiply three local terms rather than fill a full \(2\times2\times2\) table independently.

\paragraph{Computational neuroscience example.}
A graphical model for a neural experiment might include stimulus \(S\), latent attention state \(Z\), neural response \(R\), and choice \(C\). A possible DAG is
\[
S\to R,\qquad Z\to R,\qquad R\to C,\qquad S\to C.
\]
This graph expresses assumptions about which variables directly predict others. It does not prove biological causality by itself. It defines a probabilistic factorization that can be tested, fitted, criticized, and compared with alternatives.

\paragraph{AI and machine-learning example.}
Naive Bayes classification uses a label \(Y\) as a parent of features \(X_1,\ldots,X_d\):
\[
p(y,x_1,\ldots,x_d)=p(y)\prod_{j=1}^d p(x_j\mid y).
\]
The conditional-independence assumption is strong and often false literally, but it can produce useful classifiers because it drastically reduces the number of parameters. Conditional random fields, Markov random fields, and factor graphs use undirected or factorized structure for sequences, images, and structured prediction.

\paragraph{Common misconceptions and failure modes.}
First, an edge is not automatically a physical causal connection. In a probabilistic graphical model, edges encode factorization assumptions; causal interpretation requires additional intervention assumptions. Second, missing edges encode conditional independences only relative to the chosen graph and model class. Third, undirected potentials are not probabilities by themselves; they must be normalized by \(Z\). Fourth, conditioning can create or remove dependencies depending on graph structure, so intuition from marginal correlations can mislead.

\paragraph{Exercises.}
\begin{enumerate}[label=\textbf{34.1.\arabic*.}, leftmargin=*]
\item \textbf{Mechanical.} For the DAG \(A\to B\), \(A\to C\), \(B\to D\), write the factorization of \(p(a,b,c,d)\).
\item \textbf{Proof or derivation.} For \(Z\to X\) and \(Z\to Y\), prove from the factorization that \(X\perp Y\mid Z\).
\item \textbf{Numerical/computational.} Using the binary example in this section, compute \(P(Z=0,X=1,Y=1)\).
\item \textbf{Interpretation.} In a neural graphical model, why should an edge from \(R\) to \(C\) be interpreted cautiously if the data are observational?
\item \textbf{Misconception/debugging.} A student says, "If two variables are correlated, the graph must contain an edge between them." Explain why shared parents or paths through other variables can create marginal correlation.
\end{enumerate}

\subsection{34.2 Latent variable models}

\paragraph{Guiding question.}
How can hidden variables explain observed data, and how do we reverse the generative direction to infer hidden causes?

\paragraph{Beginner-facing intuition.}
A latent-variable model says that observations are generated from hidden variables. The model direction is
\[
\text{latent cause } Z \longrightarrow \text{ observation } X.
\]
Inference goes the other way:
\[
\text{observation } X=x \longrightarrow \text{ posterior over } Z.
\]
This reversal is Bayes' theorem. The latent variable may be a cluster identity, stimulus state, movement intention, low-dimensional neural state, topic, or hidden dynamical state.

\paragraph{Formal definitions.}
Let \(Z\in\mathcal Z\) be latent and \(X\in\mathcal X\) be observed. A latent-variable generative model specifies
\[
\boxed{
p_\theta(x,z)=p_\theta(z)p_\theta(x\mid z).
}
\]
The marginal likelihood of an observation is
\[
\boxed{
p_\theta(x)=\sum_{z\in\mathcal Z}p_\theta(z)p_\theta(x\mid z)
}
\]
for discrete \(Z\), or
\[
p_\theta(x)=\int p_\theta(z)p_\theta(x\mid z)\,dz
\]
for continuous \(Z\). The posterior is
\[
\boxed{
p_\theta(z\mid x)=\frac{p_\theta(z)p_\theta(x\mid z)}{p_\theta(x)}.
}
\]
The denominator is the normalization that makes posterior probabilities sum or integrate to one.

\paragraph{Core mechanism: exact posterior in a two-component mixture.}
Suppose \(Z\in\{1,2\}\) chooses one of two components:
\[
P(Z=1)=\pi,\qquad P(Z=2)=1-\pi.
\]
The observation model is \(p(x\mid Z=k)=p_k(x)\). Then
\[
P(Z=1\mid x)
=
\frac{\pi p_1(x)}{\pi p_1(x)+(1-\pi)p_2(x)}.
\]
The posterior responsibility of component \(1\) is high when its prior weight \(\pi\) is large or its likelihood \(p_1(x)\) is large relative to the other component.

\paragraph{Concrete numerical example.}
Let \(\pi=0.3\). At an observed value \(x\), suppose
\[
p_1(x)=0.8,\qquad p_2(x)=0.2.
\]
Then
\[
P(Z=1\mid x)
=\frac{0.3\cdot0.8}{0.3\cdot0.8+0.7\cdot0.2}
=\frac{0.24}{0.38}
\approx0.632.
\]
Even though component \(1\) had lower prior probability, the observation made it more plausible because it fit \(x\) much better.

\paragraph{Sequential Bayes.}
If observations arrive one at a time and a fixed latent state \(Z\) is being inferred, Bayes' theorem can be applied iteratively. After observing \(x_1\),
\[
p(z\mid x_1)\propto p(z)p(x_1\mid z).
\]
When \(x_2\) arrives, the previous posterior becomes the new prior:
\[
\boxed{
p(z\mid x_1,x_2)
\propto
p(z\mid x_1)p(x_2\mid z)
}
\]
assuming \(x_1\) and \(x_2\) are conditionally independent given \(Z\). This posterior-after-one-observation-becomes-prior-for-the-next pattern is the basis of filtering and state estimation.

\paragraph{Computational neuroscience example.}
In a simple latent-state model, \(Z\) may indicate whether an animal is attentive or inattentive, while \(X\) is a population response vector. The generative model \(p(z)p(x\mid z)\) says that the hidden state changes the distribution of neural activity. Given an observed response \(x\), posterior inference computes \(p(z\mid x)\), the model's belief about the hidden state. This is a statistical summary, not direct access to a biological variable unless validated by independent evidence.

\paragraph{AI and machine-learning example.}
Gaussian mixture models use a discrete latent cluster \(Z\) to generate a continuous observation \(X\). Topic models use a latent topic mixture to generate words. Variational autoencoders use continuous latent variables \(Z\) to generate high-dimensional observations \(X\) through a neural decoder \(p_\theta(x\mid z)\). In each case, generation runs from \(z\) to \(x\), while inference estimates \(z\) from \(x\).

\paragraph{Common misconceptions and failure modes.}
First, latent variables are model constructs; they are not automatically real hidden objects. Second, a high posterior probability can result from the prior, the likelihood, or both. Third, marginal likelihood requires summing or integrating over latent variables, which can be computationally hard. Fourth, different latent-variable models can explain the same observations, so identifiability needs attention.

\paragraph{Exercises.}
\begin{enumerate}[label=\textbf{34.2.\arabic*.}, leftmargin=*]
\item \textbf{Mechanical.} Write \(p_\theta(x)\) and \(p_\theta(z\mid x)\) for a discrete latent-variable model with \(z\in\{1,\ldots,K\}\).
\item \textbf{Proof or derivation.} Starting from \(p(x,z)=p(z)p(x\mid z)=p(x)p(z\mid x)\), derive Bayes' formula for \(p(z\mid x)\).
\item \textbf{Numerical/computational.} Let \(P(Z=1)=0.6\), \(p_1(x)=0.1\), and \(p_2(x)=0.4\). Compute \(P(Z=1\mid x)\).
\item \textbf{Interpretation.} Explain the difference between the generative direction \(p(x\mid z)\) and the inference direction \(p(z\mid x)\) in a neural latent-state model.
\item \textbf{Misconception/debugging.} A student says, "The component with the highest likelihood always has the highest posterior." Give a counterexample using unequal priors.
\end{enumerate}

\subsection{34.3 Expectation-maximization}

\paragraph{Guiding question.}
How can we estimate parameters when the likelihood depends on hidden variables we did not observe?

\paragraph{Beginner-facing intuition.}
If latent variables were observed, parameter fitting would often be easy. If parameters were known, posterior inference over latent variables would often be easy. Expectation-maximization, or EM, alternates between these two easier problems. The E-step infers a distribution over hidden variables using current parameters. The M-step updates parameters as if hidden variables were softly filled in by that posterior.

EM is not magic global optimization. It is a coordinate-ascent method on a lower bound to the log marginal likelihood.

\paragraph{Formal definitions.}
For observations \(x_1,\ldots,x_n\) and latent variables \(z_1,\ldots,z_n\), the marginal log-likelihood is
\[
\ell(\theta)=\sum_{i=1}^n \log p_\theta(x_i)
=\sum_{i=1}^n \log \sum_{z_i} p_\theta(x_i,z_i)
\]
for discrete latent variables. The sum inside the log couples inference and parameter estimation.

Given current parameters \(\theta^{\mathrm{old}}\), the E-step computes
\[
q_i(z_i)=p_{\theta^{\mathrm{old}}}(z_i\mid x_i).
\]
The M-step updates
\[
\theta^{\mathrm{new}}
\in
\operatorname*{arg\,max}_\theta
\sum_{i=1}^n
\mathbb E_{q_i}
\left[
\log p_\theta(x_i,z_i)
\right].
\]

\paragraph{Core derivation: Jensen lower bound and EM.}
For any distribution \(q(z)\),
\[
\log p_\theta(x)
=\log\sum_z q(z)\frac{p_\theta(x,z)}{q(z)}.
\]
By Jensen's inequality,
\[
\log p_\theta(x)
\geq
\sum_z q(z)\log\frac{p_\theta(x,z)}{q(z)}
=
\mathbb E_q[\log p_\theta(x,z)]
+H(q).
\]
This lower bound is often called an evidence lower bound, or ELBO:
\[
\mathcal L(q,\theta)
=\mathbb E_q[\log p_\theta(x,z)]-\mathbb E_q[\log q(z)].
\]
It also satisfies
\[
\boxed{
\log p_\theta(x)
=\mathcal L(q,\theta)
+D_{\mathrm{KL}}\bigl(q(z)\|p_\theta(z\mid x)\bigr).
}
\]
For fixed \(\theta\), the bound is tight when \(q(z)=p_\theta(z\mid x)\). EM uses this fact: the E-step sets \(q\) to the current posterior, making the bound tight at \(\theta^{\mathrm{old}}\); the M-step increases the bound with respect to \(\theta\). This guarantees nondecreasing likelihood under exact E and M steps, though not convergence to the global maximum.

\paragraph{Concrete example: mixture responsibilities.}
For a two-component mixture with parameters \(\pi_k\) and component densities \(p_k(x;\theta_k)\), the E-step computes responsibilities
\[
r_{ik}
=P(Z_i=k\mid x_i)
=
\frac{\pi_k p_k(x_i;\theta_k)}
{\sum_{\ell=1}^2 \pi_\ell p_\ell(x_i;\theta_\ell)}.
\]
The M-step for mixture weights is
\[
\pi_k^{\mathrm{new}}=\frac1n\sum_{i=1}^n r_{ik}.
\]
For Gaussian components with known variance, the mean update is a weighted average:
\[
\mu_k^{\mathrm{new}}
=
\frac{\sum_i r_{ik}x_i}{\sum_i r_{ik}}.
\]
If two observations \(x_1=0\), \(x_2=10\) have responsibilities for component 1 equal to \(0.9\) and \(0.2\), then
\[
\mu_1^{\mathrm{new}}
=\frac{0.9(0)+0.2(10)}{0.9+0.2}
=\frac{2}{1.1}
\approx1.82.
\]

\paragraph{Computational neuroscience example.}
EM can fit mixture models of neural activity states, hidden Markov models for switching neural dynamics, or latent-factor models when latent states are not observed. In a hidden Markov model, the E-step computes posterior probabilities of hidden state sequences or state occupancies, and the M-step updates transition and emission parameters. The inferred states are model-dependent summaries and should be validated against behavior, stimuli, or held-out data.

\paragraph{AI and machine-learning example.}
EM underlies Gaussian mixture clustering, missing-data models, topic models, and some speech-recognition and sequence models. Variational autoencoders can be viewed as amortized, approximate descendants of the same idea: infer latent variables under current parameters, then update generative parameters, but with a learned approximate inference network instead of exact E-steps.

\paragraph{Common misconceptions and failure modes.}
First, EM does not fill in a single best latent value; it usually uses posterior distributions or responsibilities. Second, EM can converge to local optima and depends on initialization. Third, exact E-steps may be impossible in complex models, leading to variational or Monte Carlo EM. Fourth, increasing likelihood does not guarantee that the latent variables are scientifically meaningful.

\paragraph{Exercises.}
\begin{enumerate}[label=\textbf{34.3.\arabic*.}, leftmargin=*]
\item \textbf{Mechanical.} Define the E-step and M-step for a latent-variable model in words and formulas.
\item \textbf{Proof or derivation.} Use Jensen's inequality to derive the ELBO \(\mathcal L(q,\theta)\leq\log p_\theta(x)\).
\item \textbf{Numerical/computational.} For responsibilities \(r_{i1}=(0.8,0.6,0.1)\) and observations \(x=(1,2,9)\), compute the M-step update for \(\mu_1\) in a Gaussian mixture with known variance.
\item \textbf{Interpretation.} In a neural hidden-state model, what does a responsibility \(r_{ik}=0.7\) mean?
\item \textbf{Misconception/debugging.} A student says, "EM always finds the true clusters." Explain the roles of local optima, model misspecification, and initialization.
\end{enumerate}

\subsection{34.4 Approximate inference}

\paragraph{Guiding question.}
What can we do when the exact posterior \(p(z\mid x)\) is too expensive to compute?

\paragraph{Beginner-facing intuition.}
Exact inference requires normalizing over all latent possibilities:
\[
p(z\mid x)=\frac{p(x,z)}{\sum_{z'}p(x,z')}
\]
or an integral over continuous latent variables. In high-dimensional models, this sum or integral can be enormous. Approximate inference replaces the exact posterior with something computable: a Gaussian near a mode, a simpler variational distribution, samples, messages on a graph, or particles that track a filtering distribution through time.

Different approximations make different tradeoffs between accuracy, computation, and diagnostic transparency.

\paragraph{Formal setup.}
Let \(p_\theta(x,z)\) be a joint model. The posterior is
\[
p_\theta(z\mid x)=\frac{p_\theta(x,z)}{p_\theta(x)},
\qquad
p_\theta(x)=\int p_\theta(x,z)\,dz.
\]
The hard part is often \(p_\theta(x)\), the evidence. Approximate inference aims to compute expectations such as
\[
\mathbb E[g(Z)\mid X=x]
\]
without exact normalization.

\paragraph{Core derivation: variational identity.}
For any density or mass function \(q(z)\),
\[
\log p_\theta(x)
=
\mathbb E_q[\log p_\theta(x,z)-\log q(z)]
+D_{\mathrm{KL}}\bigl(q(z)\|p_\theta(z\mid x)\bigr).
\]
Since KL divergence is nonnegative,
\[
\boxed{
\mathcal L(q,\theta)
=\mathbb E_q[\log p_\theta(x,z)-\log q(z)]
\leq \log p_\theta(x).
}
\]
Variational inference chooses \(q\) from a tractable family \(\mathcal Q\) and maximizes \(\mathcal L(q,\theta)\), equivalently minimizing
\[
D_{\mathrm{KL}}\bigl(q(z)\|p_\theta(z\mid x)\bigr)
\]
within that family. Mean-field variational inference assumes a factorization such as
\[
q(z_1,\ldots,z_m)=\prod_{j=1}^m q_j(z_j),
\]
which simplifies computation but can miss posterior correlations.

\paragraph{Main approximation families.}
\emph{Laplace approximation.} Find a posterior mode
\[
z_0=\operatorname*{arg\,max}_z \log p(z\mid x).
\]
Approximate the log posterior by a quadratic Taylor expansion near \(z_0\):
\[
\log p(z\mid x)\approx \log p(z_0\mid x)
-\frac12(z-z_0)^\top H(z-z_0),
\]
where
\[
H=-\nabla_z^2\log p(z\mid x)\big|_{z=z_0}
\]
is positive definite at a local maximum. The approximate posterior is Gaussian:
\[
q(z)=\mathcal N(z_0,H^{-1}).
\]

\emph{Monte Carlo and MCMC.} Approximate posterior expectations by samples:
\[
\mathbb E[g(Z)\mid x]\approx \frac1M\sum_{m=1}^M g(z^{(m)}).
\]
If direct sampling from \(p(z\mid x)\) is impossible, Markov chain Monte Carlo constructs dependent samples whose stationary distribution is the posterior.

\emph{Message passing.} On a factor graph with variables \(x_i\) and factors \(\psi_a(x_{\partial a})\), suppose the joint distribution factors as
\[
p(x_1,\ldots,x_m)=\frac{1}{Z}\prod_{a\in\mathcal F}\psi_a(x_{\partial a}),
\]
where \(\partial a\) is the set of variables touching factor \(a\), \(N(i)\) is the set of factors touching variable \(i\), and \(Z\) is the normalizing constant. Sum-product passes two local message types. The variable-to-factor message is
\[
\boxed{
m_{i\to a}(x_i)=\prod_{b\in N(i)\setminus\{a\}}m_{b\to i}(x_i),
}
\]
and the factor-to-variable message is
\[
\boxed{
m_{a\to i}(x_i)
=
\sum_{x_{\partial a\setminus\{i\}}}
\psi_a(x_{\partial a})
\prod_{j\in\partial a\setminus\{i\}}m_{j\to a}(x_j).
}
\]
For continuous variables, the sum is replaced by an integral. After messages into variable \(i\) have arrived, its local belief is
\[
b_i(x_i)\propto \prod_{a\in N(i)}m_{a\to i}(x_i),
\]
with a final normalization over \(x_i\). On tree-structured factor graphs, these local sums and products compute exact marginals after a finite inward-outward pass. On loopy graphs, running the same updates repeatedly is generally approximate and may fail to converge.

\emph{Particle filtering.} For a state-space model with latent state \(Z_t\) and observation \(X_t\), filtering uses the prediction-update recursion
\[
p(z_t\mid x_{1:t-1})
=\int p(z_t\mid z_{t-1})p(z_{t-1}\mid x_{1:t-1})\,dz_{t-1},
\]
\[
\boxed{
p(z_t\mid x_{1:t})
\propto
p(x_t\mid z_t)p(z_t\mid x_{1:t-1}).
}
\]
Particles approximate the filtering distribution by weighted samples. The posterior after time \(t-1\) becomes the prior for time \(t\) after propagation through the dynamics.

\paragraph{Concrete example.}
Suppose \(Z\) is one-dimensional and the log posterior near its mode is
\[
\log p(z\mid x)\approx C-2(z-3)^2.
\]
Match this to a Gaussian log density
\[
C'-\frac12\frac{(z-\mu)^2}{\sigma^2}.
\]
Thus \(\mu=3\) and
\[
\frac{1}{2\sigma^2}=2
\quad\Longrightarrow\quad
\sigma^2=\frac14.
\]
The Laplace approximation is \(q(z)=\mathcal N(3,0.25)\).

\paragraph{Computational neuroscience example.}
In neural state-space models, latent variables may represent position, movement intention, arousal, or low-dimensional population state. Exact filtering over continuous nonlinear states is rarely available. Particle filters approximate the posterior over the current state from neural observations and dynamics. Variational methods approximate latent trajectories in larger models. Message passing is exact for some linear-Gaussian or tree-like discrete models and approximate for more complex recurrent circuits.

\paragraph{AI and machine-learning example.}
Variational inference is central in VAEs and probabilistic topic models. Laplace approximations give local Gaussian uncertainty around MAP estimates. MCMC is used when asymptotic exactness is more important than speed. Particle filters appear in tracking, robotics, sequential Monte Carlo, and differentiable filtering models. Message passing supports error-correcting codes, graphical models, and some structured prediction systems.

\paragraph{Common misconceptions and failure modes.}
First, approximate inference is not one method; each approximation has assumptions. Second, a variational posterior can be overconfident if the factorized family cannot represent correlations. Third, MCMC samples are not independent by default and require convergence diagnostics. Fourth, particle filters can degenerate when most weights collapse onto a few particles. Fifth, exact-looking formulas in a misspecified model still produce model-based posteriors, not guaranteed truths.

\paragraph{Exercises.}
\begin{enumerate}[label=\textbf{34.4.\arabic*.}, leftmargin=*]
\item \textbf{Mechanical.} List Laplace approximation, variational inference, Monte Carlo, message passing, and particle filtering, giving one sentence for what each approximates.
\item \textbf{Proof or derivation.} Derive the variational identity \(\log p(x)=\mathcal L(q,\theta)+D_{\mathrm{KL}}(q\|p(z\mid x))\).
\item \textbf{Numerical/computational.} If \(\log p(z\mid x)\approx C-\frac{1}{8}(z-5)^2\), compute the Laplace Gaussian mean and variance.
\item \textbf{Interpretation.} In particle filtering, explain the prediction step and the update step in terms of latent neural state and new neural observation.
\item \textbf{Misconception/debugging.} A student says, "Variational inference gives the posterior." Correct the statement by explaining the role of the variational family \(\mathcal Q\).
\end{enumerate}

\paragraph{Closing bridge.}
Graphical models and latent-variable inference complete the first machine-learning arc of the course. We can now define prediction functions, regularize them, place distributions on functions, and infer hidden causes behind observations. The next chapters use these tools for unsupervised representation learning, reinforcement learning, deep networks, generative modeling, neural data analysis, and theories of perception as sequential model-based inference.

\clearpage
